\section{Strategien zur nutzerzentrierten Informationsarchitektur}
\label{sec:07-01_ia-strategy-recommendations}

Auf Basis der in \textbf{\hyperref[ch:05_analysis-and-results]{Kapitel~5}} und \textbf{\hyperref[sec:06-03_organizational-ux-implications]{Abschnitt~6.3}} gewonnenen Erkenntnisse lassen sich konkrete Handlungsempfehlungen für die strategische Gestaltung einer nutzerzentrierten \Gls{ia} in organisationalen Wissenssystemen ableiten. Ziel ist es, die identifizierten Erfolgsfaktoren systematisch zu verstetigen und die erkannten Schwachstellen nachhaltig zu adressieren.

Zentral ist die institutionelle Verankerung von \Gls{ia} als kontinuierliche Gestaltungsaufgabe. Anstelle einer projektbezogenen Einmalgestaltung sollte die Systemstruktur regelmäßig überprüft und entlang realer Nutzungsmuster weiterentwickelt werden. Wiederkehrende Reviews von Navigationslogiken, Kategoriestrukturen und Benennungssystemen vermeiden strukturelle Inkonsistenzen und sichern Erwartungskonformität langfristig \cite{303:Rosenfeld2015,108:ArangoFirstPrinciples2024}. Ergänzend ist eine klare Zuweisung von Verantwortlichkeiten erforderlich, die sich auch auf einzelne Datenquellen und Inhaltsbereiche erstreckt. Jede Informationsquelle sollte über eine eindeutig benannte pflegende Ansprechperson verfügen, die Aktualität, Konsistenz und Qualität sicherstellt und inkonsistente Parallelentwicklungen verhindert \cite{309:NNGContentGovernance}.

Die empirischen Befunde unterstreichen die zentrale Rolle mentaler Modelle für Orientierung und Auffindbarkeit. Qualitative Methoden wie Interviews, Card-Sorting oder strukturierte Feedbackformate ermöglichen es, \gls{mentalemodelle} der Zielgruppen systematisch zu erfassen und in Strukturentscheidungen zu überführen \cite{905:Norman1983,303:Rosenfeld2015}. Eine enge Verzahnung von Nutzerverständnis und \Gls{ia}-Design trägt dazu bei, Navigationspfade intuitiv erschließbar zu halten, Suchaufwände zu reduzieren und die Struktur entlang realer Arbeitspraktiken auszurichten.

Einstiegs- und Übersichtsseiten spielen dabei eine wichtige Rolle, da sie als zentrale Orientierungspunkte für die Nutzung dienen. Durch die klare Priorisierung häufig benötigter Inhalte, eine verständliche visuelle Struktur und konsistente Verlinkungen in tieferliegende Bereiche unterstützen sie sowohl den schnellen Überblick als auch das gezielte Auffinden von Informationen \cite{113:NNGIntranetGuidelines,307:DWGFindability}. Eine regelmäßige Überprüfung anhand realer Nutzungsszenarien stellt sicher, dass diese Seiten auch bei veränderten Arbeitsprozessen relevant bleiben.

Gleichzeitig zeigen die Ergebnisse, dass einzelne lokale Strukturdefizite die Gesamtwirkung einer ansonsten konsistenten \Gls{ia} deutlich beeinträchtigen können. Kritische Inhaltsbereiche sollten daher gezielt über \Gls{usability}-Tests überprüft und identifizierte Schwachstellen systematisch priorisiert werden \cite{121:NNGUsabilityTesting}. Bereits kleinere strukturelle Anpassungen, etwa durch präzisere Benennungen, zusätzliche Querverlinkungen oder alternative Einstiegspfade, können substanzielle Verbesserungen der Nutzungseffizienz bewirken \cite{303:Rosenfeld2015}.

Suchfunktion und \Gls{navigationsstruktur} sollten zudem nicht isoliert optimiert, sondern als komplementäre Zugangssysteme verstanden werden. Verbesserungen in Metadatenqualität, Filterlogik und Ranking erhöhen die Wirksamkeit der Suche, während konsistente Navigationspfade die Abhängigkeit von Suchprozessen reduzieren \cite{111:NNGSearchVisible,309:NNGContentGovernance}. Eine integrierte Weiterentwicklung beider Systeme erhöht die Robustheit der Informationsfindung über unterschiedliche Nutzungskontexte hinweg.

Die in dieser Arbeit eingesetzte Kombination aus quantitativen Nutzungsmetriken und qualitativen Rückmeldungen verdeutlicht schließlich die Bedeutung kontinuierlicher Evaluation. Iterative Validierungszyklen ermöglichen eine evidenzbasierte Weiterentwicklung der Struktur und unterstützen eine lernorientierte Organisationspraxis \cite{124:Triangulation2023,121:NNGUsabilityTesting}. \Gls{leanux}-Ansätze fördern hierbei kurze Feedbackschleifen und eine adaptive Weiterentwicklung, wodurch strukturelle Fehlentwicklungen frühzeitig erkannt und korrigiert werden können \cite{424:GothelfSeiden2016}.

Insgesamt zeigen die Empfehlungen, dass eine nutzerzentrierte \Gls{ia} weniger durch statische Regelwerke als durch kontinuierliche Lern- und Anpassungsprozesse entsteht. Die strategische Verankerung von Strukturpflege, Nutzerfeedback und iterativer Evaluation bildet die Grundlage für langfristige Gebrauchstauglichkeit und nachhaltige Akzeptanz organisationaler Wissensplattformen.
